\begin{definition}
Let $\Omega\subset\mathbb{C}$ be open and let $1\le K<\infty$.  A mapping
$f\colon\mathbb{C}\to\mathbb{C}$ is \emph{$K$-quasiregular on $\Omega$} if
\begin{enumerate}
\item[(a)] $f\in W^{1,2}_{\mathrm{loc}}(\Omega)$, in the sense of
\noderef{MemWLocOn}; and
\item[(b)] there are weak Wirtinger derivatives $a=\partial f/\partial z$ and
$b=\partial f/\partial\bar z$ of $f$ on $\Omega$, in the sense of
\noderef{IsWeakWirtingerOn}, such that
\[
  \bigl|b(z)\bigr|\le\frac{K-1}{K+1}\bigl|a(z)\bigr|
\]
for almost every $z\in\Omega$, ``almost every'' referring to two-dimensional
Lebesgue measure.
\end{enumerate}

This is the definition of $K$-quasiregularity given at paper.tex:129--133,
together with its differential inequality
\[
  \Bigl|\frac{\partial f(z)}{\partial\bar z}\Bigr|
  \le\frac{K-1}{K+1}\Bigl|\frac{\partial f(z)}{\partial z}\Bigr|
  \tag{11}
\]
for almost every $z\in\Omega$.

Two remarks on the reading of (11).

First, the derivatives in (b) are the weak ones.  This is what the paper's
sentence says: the only hypothesis $f$ carries where (11) is stated is
$f\in W^{1,2}_{\mathrm{loc}}(\Omega)$, and the two Wirtinger operators have
just been introduced, two lines earlier, as the combinations
$\frac12(\partial_x\mp i\partial_y)$; so the derivatives in (11) are the
distributional ones whose existence the Sobolev membership asserts.  Reading
them instead through the real Fr\'echet derivative of the given representative
would not formalize (11) at all.  Clause (a) constrains $f$ only up to null
sets --- by \noderef{MemWLocOn} it is a family of integral identities and
local $L^2$-conditions --- so a class in $W^{1,2}_{\mathrm{loc}}(\Omega)$ has
representatives that are differentiable at no point whatever; for such a
representative every pointwise Wirtinger derivative is $0$, by the convention
that the differential is $0$ where it does not exist, and the pointwise form
of (11) would hold vacuously however large the class's genuine distortion.
With (b) as stated this cannot happen: the weak derivatives depend only on the
class of $f$, so (11) is a condition on the class, as the paper intends.

Second, the witnessing pair in (b) is immaterial.  By the last paragraph of
\noderef{IsWeakWirtingerOn}, two pairs of weak Wirtinger derivatives of $f$ on
an open $\Omega$ agree almost everywhere there, so the validity of (b) does not
depend on which pair is chosen; and (b) for one pair gives (b) for every pair.
\end{definition}
